\section{引言}
接触广泛存在于机器人、机构、车辆、装配与颗粒系统等多体动力学问题中，其建模质量直接影响系统的受力分析、运动预测与数值稳定性。长期以来，多体动力学中的接触大多采用两类基本路线：一类是基于非穿透约束的互补类硬接触模型，另一类是基于局部压入的罚函数或顺应接触模型。前者能够较直接地表达理想刚体接触的几何约束，后者则在工程实现、数值连续性和复杂场景适应性方面更具灵活性。然而，随着复杂 CAD 几何、非凸刚体和面接触问题在仿真中的不断增加，传统接触模型的局限逐渐显现。

现有点接触模型通常将接触作用简化为单个或少量离散接触点。在球—平面、简单凸体碰撞等场景中，这种近似往往已足够有效；但在盒体面接触、边缘滚翻、足底支撑、夹持及非凸结构接触等问题中，真实接触往往发生在有限区域内，净接触力与力矩取决于接触斑的位置、范围及其上的压力分布，而不仅仅取决于若干离散点的位置。因此，点接触模型在这类场景中容易出现接触中心漂移、接触力矩估计不稳定、接触切换抖动以及对几何离散敏感等问题。换言之，复杂刚体接触的关键已不再只是“是否发生接触”，而是“接触发生在何处、以何种空间分布形式作用于系统”。

从更广义的接触建模角度看，互补类方法更适合描述宏观刚体约束，但难以自然表达有限接触斑上的分布压力及其引起的合力矩；传统顺应接触方法虽然更适于工程仿真，却往往仍以点级接触力为基本单元，尚未从根本上解决复杂面接触的空间分布表达问题。因此，在多体动力学中仍缺乏一种同时兼顾复杂几何表示能力、接触区域描述能力和工程可实现性的轻量化分布式接触模型。

符号距离场（Signed Distance Field, SDF）为这一问题提供了新的几何基础。SDF 能够统一表示任意复杂、非凸甚至带孔洞的几何形状，并可直接提供局部距离、法向、压入量及局部几何质量信息。近年来，SDF 已广泛用于碰撞检测、最小距离查询、隐式几何建模和可微接触中。然而，在多数现有做法中，SDF 仍主要充当几何查询工具，即用于获取最近点、法向或穿透深度，随后仍通过单点罚函数构造接触力。这样虽然提升了复杂几何的处理能力，却未能充分利用 SDF 在接触区域生成与空间分布建模方面的潜力。

基于上述认识，本文关注如下问题：能否以 SDF 为统一几何表示，在多体动力学中直接构造复杂刚体之间的接触斑与压力场，从而获得更合理的分布式接触力与力矩？围绕这一问题，本文提出一种基于 SDF 的接触斑—压力场建模方法。该方法不再将接触作用收缩为若干离散点力，而是首先基于双边 SDF 提取候选接触斑，再在接触斑上定义局部压力密度，并通过面分布积分得到净接触 wrench。该方法可视为介于传统点接触模型与完整连续体接触模型之间的一种中尺度接触建模框架：一方面，它保留多体动力学中顺应接触的轻量化求解优势；另一方面，它引入接触区域和分布压力的概念，以增强复杂面接触和接触力矩的表达能力。

本文的目标并非提出一个停留在概念层面的几何模型，而是构建一种能够嵌入一般多体动力学求解流程中的复杂接触建模方法，并通过系统的基准算例验证其在接触斑表达、净力矩预测和动态响应连续性方面的有效性。

本文的主要贡献包括：
\begin{enumerate}[label=(\arabic*)]
    \item 提出一种基于双边 SDF、一致性筛选和面元离散的复杂刚体接触斑构造方法；
    \item 建立由局部压入驱动的分布式压力场顺应接触模型，以及相应的净接触力与净接触力矩积分框架；
    \item 给出与多体动力学求解流程耦合的数值算法与实现要点；
    \item 设计面向复杂面接触、边角接触和非凸几何接触的基准验证方案，用于与传统点级接触模型进行系统比较。
\end{enumerate}
